\section{Verification experiments and implementation traceability}\label{sec:C25:experiments}
A verification study should predeclare the property, candidate population, observation interface, and adjudication process. It should also distinguish discovery from confirmation. Cases used to design a new gate belong to development; untouched cases test whether the gate generalizes. Repeatedly inspecting failures and adjusting the gate against the same holdout turns the holdout into a training resource even if no model weights change.

The implementation ledger links a source claim to an equation or algorithm, a versioned code path, configuration, dataset split, and observed result. For ROME, inspect the intervention and editing implementation separately; a successful edit does not validate every localization claim. For sparse autoencoders, inspect training, feature extraction, and downstream metrics separately. For executable witnesses, preserve the native test environment and the requirement argument. These boundaries prevent one successful component from lending unwarranted credibility to another.

Negative controls are particularly valuable. Patch a random activation site, use an irrelevant feature description, supply a test that distinguishes an incidental reference choice, or give a judge a trajectory for a related but incompatible instruction. A valid verification method should reject or limit these controls appropriately. Failing a negative control identifies a concrete weakness in causal interpretation or specification handling.

Uncertainty must follow the sampling unit. Multiple prompts for one fact, several features from one model layer, and many candidate patches for one task are not independent replications of the entire method. Use task-, fact-, model-, or repository-level grouping as appropriate. Report heterogeneity rather than only an aggregate mean, especially when rare errors have disproportionate consequences.

The practical implication is that verification is an evidence-producing system with its own model, tools, data, and failure modes. It should be versioned, monitored, and independently evaluated. Internal inspection can reveal useful mechanisms; executable checks can falsify artifacts; structured evidence can expose unsupported claims. Their combined strength comes from complementary observation and intervention, not from relabeling a second model as an infallible critic.
